Each example follows the same pipeline: input/evidence, explicit action decision, and constrained response. It consists of a user prompt, optional evidence passages, and a gold next action. The model must output an action and a concise response. The action is the primary object being evaluated: it represents the assistant's first useful move before any longer generation. The four actions are:

\begin{itemize}
    \item \texttt{answer}: the request is sufficiently specified, the premise is acceptable, and the evidence supports a concrete answer.
    \item \texttt{ask}: the request is missing information that the user can plausibly provide in one short follow-up.
    \item \texttt{challenge}: the prompt contains a false or stale premise that should be corrected before answering.
    \item \texttt{abstain}: reliable support is missing or irreconcilably conflicting, and a follow-up would not resolve the issue.
\end{itemize}

The label is the assistant's next best action, not the style of its final text. This distinction separates the task from both standard question answering and generic refusal detection\citep{min-etal-2020-ambigqa,rajpurkar-etal-2018-know,geifman-el-yaniv-2017-selective}. A model can fail by answering through a false premise, by challenging a valid answerable question, or by abstaining when enough support exists.

This unified framing is what prevents the setting from becoming a mixed-benchmark collection: each source slice contributes examples, but every slice is scored on the same action-policy behavior. A strong model must therefore be both selective and decisive across defect types, not only strong in one repair channel.

We apply a fixed decision order when assigning gold actions. If the request is missing essential user-provided information that one short follow-up could resolve, the gold action is \texttt{ask}. If the request is specified but contains a false or stale premise, the gold action is \texttt{challenge}. If the premise is acceptable and the evidence supports a concrete answer, the gold action is \texttt{answer}. Otherwise, the gold action is \texttt{abstain}. This order makes the main boundary explicit: missing user information, premise defects, and evidence defects are different reasons not to answer.

We use action accuracy as the primary discrete metric and report per-slice scores to avoid hiding structured failures. We also report over-answer rate, defined as the fraction of non-\texttt{answer} gold examples where the model predicts \texttt{answer}. Finally, we use an average utility score as a diagnostic summary. For answerable examples, a correct direct answer receives $1.0$; an incorrect direct answer or unsupported abstention receives $-0.5$; unnecessary clarification receives $-0.25$; and wrongly challenging a valid prompt receives $-0.75$. For non-answer examples, the correct repair action receives $0.5$, direct over-answering receives $-1.0$, and other non-answer mistakes receive $-0.25$. These weights encode an asymmetric harm ordering, not a universal user-cost model. We therefore never interpret utility alone: tables pair it with action accuracy, over-answer rate, answer containment, JSON parse rate, and per-slice metrics, following the broader lesson that confidence and uncertainty summaries must be checked against behavior rather than treated as self-validating scores\citep{guo-etal-2017-calibration,kadavath-etal-2022-language,kuhn-etal-2023-semantic}.

For this paper-facing stage, we report all action types in the same table/figure layout, but interpret non-answer classes with special care: \texttt{ask} and \texttt{abstain} are paper-relevant behaviors even when some future slices are not yet scaled to full size. We explicitly separate evidence claims by slice and action to prevent a single scalar score from hiding boundary mistakes.
